\documentclass[12pt,a4paper]{article}

% ── Packages ──────────────────────────────────────────────
\usepackage[margin=2.5cm]{geometry}
\usepackage{graphicx}
\usepackage{booktabs}
\usepackage{array}
\usepackage{xcolor}
\usepackage{titlesec}
\usepackage{fancyhdr}
\usepackage{amsmath}
\usepackage{hyperref}
\usepackage{enumitem}
\usepackage{caption}
\usepackage{subcaption}
\usepackage{float}
\usepackage{multirow}
\usepackage{colortbl}

% ── Colors ────────────────────────────────────────────────
\definecolor{suslblue}{RGB}{0,70,127}
\definecolor{lightgray}{RGB}{245,245,245}
\definecolor{darkgray}{RGB}{80,80,80}
\definecolor{successgreen}{RGB}{39,174,96}
\definecolor{alertred}{RGB}{192,57,43}

% ── Section Formatting ────────────────────────────────────
\titleformat{\section}{\large\bfseries\color{suslblue}}{}{0em}{}[\titlerule]
\titleformat{\subsection}{\normalsize\bfseries\color{darkgray}}{}{0em}{}

% ── Header / Footer ───────────────────────────────────────
\pagestyle{fancy}
\fancyhf{}
\rhead{\textcolor{darkgray}{\small SE 8101 Research Project}}
\lhead{\textcolor{darkgray}{\small S.Shanthosh | 20APSE4882}}
\rfoot{\textcolor{darkgray}{\small Page \thepage}}
\lfoot{\textcolor{darkgray}{\small Sabaragamuwa University of Sri Lanka}}

% ── Hyperlink Setup ───────────────────────────────────────
\hypersetup{colorlinks=true, linkcolor=suslblue, urlcolor=suslblue}

% ═══════════════════════════════════════════════════════════
\begin{document}

% ── Title Page ────────────────────────────────────────────
\begin{titlepage}
  \centering
  \vspace*{1cm}
  {\Large\bfseries SABARAGAMUWA UNIVERSITY OF SRI LANKA\\[0.3cm]}
  {\large Faculty of Computing\\
          Department of Software Engineering\\[0.5cm]}
  {\normalsize BSc Honours Degree Programme in Software Engineering\\[2cm]}

  \rule{\linewidth}{1.5pt}\\[0.4cm]
  {\LARGE\bfseries\color{suslblue}
    Implementation Progress Report\\[0.2cm]}
  {\large SE 8101 --- Research Project\\[0.3cm]}
  \rule{\linewidth}{1.5pt}\\[1.5cm]

  \begin{tabular}{ll}
    \textbf{Student Name}   & S.Shanthosh \\[0.3cm]
    \textbf{Index Number}   & 20APSE4882  \\[0.3cm]
    \textbf{Supervisor}     & Mrs.\ WMLS Abeythunga \\[0.3cm]
    \textbf{Date}           & 27th May 2026 \\[0.3cm]
    \textbf{Report Type}    & Daily Implementation Update \\
  \end{tabular}\\[2cm]

  \colorbox{lightgray}{\parbox{0.9\linewidth}{
    \centering\vspace{0.3cm}
    {\bfseries Research Topic}\\[0.2cm]
    \textit{An Empirical Evaluation of Deep Learning Models for Enhanced and
    Scalable Bug Detection and Classification in Large-Scale Software Systems}
    \vspace{0.3cm}
  }}

  \vfill
  {\small Prepared on: 27 May 2026}
\end{titlepage}

% ── Table of Contents ─────────────────────────────────────
\tableofcontents
\newpage

% ═══════════════════════════════════════════════════════════
\section{Overview of Today's Progress}

This report documents the implementation work completed on 27th May 2026,
covering four major phases: environment configuration, exploratory data
analysis, data preprocessing, and model training. A total of six models
were trained and evaluated during this session — four traditional machine
learning models and two deep learning variants.

\begin{table}[H]
  \centering
  \caption{Implementation Progress Summary}
  \begin{tabular}{clcc}
    \toprule
    \textbf{Phase} & \textbf{Task} & \textbf{Status} & \textbf{Output} \\
    \midrule
    1 & Environment setup \& library installation & \textcolor{successgreen}{\textbf{Done}} & All libraries ready \\
    2 & Dataset acquisition (JM1 -- PROMISE)      & \textcolor{successgreen}{\textbf{Done}} & 10,885 samples \\
    3 & Exploratory Data Analysis (EDA)           & \textcolor{successgreen}{\textbf{Done}} & 4 charts saved \\
    4 & Data Preprocessing pipeline               & \textcolor{successgreen}{\textbf{Done}} & 6 split files saved \\
    5 & Baseline ML models (RF + NB)              & \textcolor{successgreen}{\textbf{Done}} & Best F1: 0.4265 \\
    6 & CNN-LSTM Deep Learning model              & \textcolor{successgreen}{\textbf{Done}} & Best Recall: 0.755 \\
    \bottomrule
  \end{tabular}
\end{table}

% ═══════════════════════════════════════════════════════════
\section{Dataset Description}

\subsection{JM1 Dataset (PROMISE Repository)}

The JM1 dataset, sourced from the PROMISE software engineering repository,
was selected as the primary dataset for this phase of the research.
It is a NASA software defect dataset containing software modules written in C.

\begin{table}[H]
  \centering
  \caption{JM1 Dataset Characteristics}
  \begin{tabular}{ll}
    \toprule
    \textbf{Property}       & \textbf{Value} \\
    \midrule
    Total software modules  & 10,885 \\
    Number of features      & 21 software complexity metrics \\
    Defective modules       & 2,106 (19.3\%) \\
    Non-defective modules   & 8,779 (80.7\%) \\
    Class imbalance ratio   & 4.2 : 1 \\
    Missing values          & 25 (in 5 columns) \\
    Feature types           & All numeric \\
    \bottomrule
  \end{tabular}
\end{table}

\subsection{Feature Description}

The 21 features comprise Halstead complexity metrics and McCabe's
cyclomatic complexity measures, including:

\begin{itemize}[leftmargin=1.5cm]
  \item \textbf{loc} -- Lines of Code
  \item \textbf{v(g)} -- Cyclomatic Complexity
  \item \textbf{ev(g), iv(g)} -- Essential and Design Complexity
  \item \textbf{n, v, l, d, i, e, b, t} -- Halstead suite of metrics
  \item \textbf{lOCode, lOComment, lOBlank} -- Source line counts
  \item \textbf{uniq\_Op, uniq\_Opnd, total\_Op, total\_Opnd} -- Operator/operand counts
  \item \textbf{branchCount} -- Number of branching points
\end{itemize}

% ═══════════════════════════════════════════════════════════
\section{Exploratory Data Analysis}

\subsection{Class Imbalance Analysis}

The EDA confirmed a significant class imbalance in the dataset, with
defective modules representing only 19.3\% of all samples. This is
consistent with real-world software systems where buggy modules are
inherently rare. A naive classifier that always predicts ``non-defective''
would achieve 80.7\% accuracy while catching zero bugs, confirming
that \textbf{F1-Score and Recall} must be used as the primary evaluation
metrics in this study.

\subsection{Feature Analysis}

Analysis of mean feature values per class revealed that \textbf{all 20 of the
21 complexity features are significantly elevated in defective modules}.
Key findings include:

\begin{table}[H]
  \centering
  \caption{Top Features Distinguishing Buggy vs Non-Buggy Modules}
  \begin{tabular}{lrrl}
    \toprule
    \textbf{Feature} & \textbf{Buggy Mean} & \textbf{Clean Mean} & \textbf{Ratio} \\
    \midrule
    t (Halstead Time)      & 6,285.2  & 1,029.6 & 6.1$\times$ \\
    e (Halstead Effort)    & 113,134  & 18,533  & 6.1$\times$ \\
    v (Halstead Volume)    & 1,422.4  & 494.2   & 2.9$\times$ \\
    loc (Lines of Code)    & 80.4     & 32.8    & 2.4$\times$ \\
    branchCount            & 21.6     & 8.8     & 2.4$\times$ \\
    v(g) (Cyclomatic)      & 11.9     & 5.0     & 2.4$\times$ \\
    \bottomrule
  \end{tabular}
\end{table}

The correlation analysis identified \textit{loc}, \textit{branchCount},
\textit{v(g)}, and \textit{n} as the features most correlated with the
defect label.

% ═══════════════════════════════════════════════════════════
\section{Data Preprocessing Pipeline}

The following preprocessing pipeline was applied sequentially:

\begin{enumerate}[leftmargin=1.5cm]
  \item \textbf{Missing Value Imputation:} 25 missing values (across
    \textit{uniq\_Op}, \textit{uniq\_Opnd}, \textit{total\_Op},
    \textit{total\_Opnd}, \textit{branchCount}) were filled using
    column-wise median imputation, which is robust against skewed
    distributions common in software metrics.

  \item \textbf{Outlier Handling:} Winsorization was applied at the
    99th percentile for each feature to cap extreme values without
    removing any data rows.

  \item \textbf{Stratified Train/Test Split:} The dataset was split
    into 80\% training (8,708 samples) and 20\% test (2,177 samples)
    using stratified sampling (\texttt{random\_state=42}) to maintain
    the original class ratio in both subsets.

  \item \textbf{Min-Max Normalization:} All features were scaled to
    the range [0, 1]. The scaler was fitted exclusively on the training
    set to prevent data leakage into the test set.

  \item \textbf{BorderlineSMOTE:} Applied only to the training set to
    generate 5,338 synthetic minority-class (buggy) samples, achieving
    a perfectly balanced 1:1 ratio (7,023 each class). The test set
    was deliberately left at its natural imbalanced ratio to reflect
    real-world conditions.
\end{enumerate}

\begin{table}[H]
  \centering
  \caption{Dataset Sizes After Preprocessing}
  \begin{tabular}{lrrr}
    \toprule
    \textbf{Split} & \textbf{Total} & \textbf{Buggy} & \textbf{Clean} \\
    \midrule
    Train (raw)    & 8,708  & 1,685 & 7,023 \\
    Train (SMOTE)  & 14,046 & 7,023 & 7,023 \\
    Test (real)    & 2,177  & 421   & 1,756 \\
    \bottomrule
  \end{tabular}
\end{table}

% ═══════════════════════════════════════════════════════════
\section{Baseline Machine Learning Models}

\subsection{Models Implemented}

Two traditional ML classifiers were trained as baselines, each under
two conditions: without SMOTE (imbalanced training data) and with SMOTE
(balanced training data).

\begin{enumerate}[leftmargin=1.5cm]
  \item \textbf{Random Forest} -- 200 estimators, max depth 15,
    \texttt{random\_state=42}
  \item \textbf{Naïve Bayes} -- Gaussian NB with default parameters
\end{enumerate}

\subsection{Baseline Results}

\begin{table}[H]
  \centering
  \caption{Baseline Model Performance on Test Set}
  \begin{tabular}{lcccccc}
    \toprule
    \textbf{Model} & \textbf{Accuracy} & \textbf{Precision}
      & \textbf{Recall} & \textbf{F1-Score} & \textbf{AUC} \\
    \midrule
    RF (no SMOTE)  & 0.8107 & 0.5280 & 0.2019 & 0.2921 & 0.7365 \\
    \rowcolor{lightgray}
    RF (SMOTE)     & 0.7653 & 0.4043 & 0.4513 & \textbf{0.4265} & 0.7350 \\
    NB (no SMOTE)  & 0.7901 & 0.4412 & 0.3207 & 0.3714 & 0.6800 \\
    NB (SMOTE)     & 0.7735 & 0.4037 & 0.3587 & 0.3799 & 0.6792 \\
    \bottomrule
  \end{tabular}
\end{table}

\subsection{Key Observations}

\begin{itemize}[leftmargin=1.5cm]
  \item Random Forest with SMOTE achieved the highest F1-Score of
    \textbf{0.4265}, establishing the baseline target to beat.
  \item Without SMOTE, RF achieved 81.1\% accuracy but caught only
    \textbf{20.2\%} of actual bugs (85/421) -- demonstrating the
    accuracy paradox in imbalanced classification.
  \item SMOTE improved Recall from 0.202 to 0.451 in Random Forest
    (\textbf{+123\%}), confirming its effectiveness.
  \item Random Forest significantly outperformed Naïve Bayes across
    all metrics, consistent with existing literature.
\end{itemize}

% ═══════════════════════════════════════════════════════════
\section{CNN-LSTM Deep Learning Model}

\subsection{Architecture}

A hybrid CNN-LSTM architecture was designed and implemented in PyTorch
to leverage both local feature pattern detection (CNN) and sequential
dependency modelling (LSTM) within the 21 software metrics.

\begin{table}[H]
  \centering
  \caption{CNN-LSTM Architecture Summary}
  \begin{tabular}{lll}
    \toprule
    \textbf{Layer} & \textbf{Configuration} & \textbf{Purpose} \\
    \midrule
    Input         & (batch, 21, 1)              & 21 metric features \\
    Conv1D        & 64 filters, kernel=3        & Local pattern detection \\
    BatchNorm     & --                          & Training stabilisation \\
    Conv1D        & 128 filters, kernel=3       & Deep feature extraction \\
    BatchNorm     & --                          & Training stabilisation \\
    LSTM          & hidden=64, layers=1         & Sequential dependencies \\
    Dropout       & rate=0.4                    & Regularisation \\
    Dense         & 32 units, ReLU              & Feature compression \\
    Dense         & 1 unit, Sigmoid             & Bug probability output \\
    \midrule
    \multicolumn{2}{l}{\textbf{Total Trainable Parameters}} & 77,121 \\
    \bottomrule
  \end{tabular}
\end{table}

\subsection{Focal Loss -- Research Contribution}

Two loss functions were evaluated:

\begin{enumerate}[leftmargin=1.5cm]
  \item \textbf{Binary Cross-Entropy (BCE):} Standard loss function
    used as the deep learning baseline.
  \item \textbf{Focal Loss} (research contribution): Originally
    proposed by Lin et al.\ (2017) for object detection, adapted here
    for software defect prediction. The loss is defined as:

    \begin{equation}
      \text{FL}(p_t) = -\alpha (1 - p_t)^{\gamma} \log(p_t)
    \end{equation}

    where $\alpha = 0.75$ weights the minority (buggy) class, and
    $\gamma = 2.0$ down-weights easy examples so training focuses on
    hard-to-classify defective modules.
\end{enumerate}

\subsection{CNN-LSTM Results}

\begin{table}[H]
  \centering
  \caption{CNN-LSTM Model Performance on Test Set (30 epochs)}
  \begin{tabular}{lccccc}
    \toprule
    \textbf{Model} & \textbf{Accuracy} & \textbf{Precision}
      & \textbf{Recall} & \textbf{F1-Score} & \textbf{AUC} \\
    \midrule
    CNN-LSTM (BCE)   & 0.6647 & 0.3057 & 0.5772 & 0.3997 & 0.6808 \\
    \rowcolor{lightgray}
    CNN-LSTM (Focal) & 0.5581 & 0.2702 & \textbf{0.7553} & 0.3980 & 0.7001 \\
    \bottomrule
  \end{tabular}
\end{table}

\subsection{Comparative Analysis}

\begin{table}[H]
  \centering
  \caption{Full Comparison -- All Models on Test Set}
  \begin{tabular}{lccccc}
    \toprule
    \textbf{Model} & \textbf{Accuracy} & \textbf{Precision}
      & \textbf{Recall} & \textbf{F1} & \textbf{AUC} \\
    \midrule
    RF (no SMOTE)          & 0.8107 & 0.5280 & 0.2019 & 0.2921 & 0.7365 \\
    RF (SMOTE)             & 0.7653 & 0.4043 & 0.4513 & 0.4265 & 0.7350 \\
    NB (no SMOTE)          & 0.7901 & 0.4412 & 0.3207 & 0.3714 & 0.6800 \\
    NB (SMOTE)             & 0.7735 & 0.4037 & 0.3587 & 0.3799 & 0.6792 \\
    CNN-LSTM (BCE)         & 0.6647 & 0.3057 & 0.5772 & 0.3997 & 0.6808 \\
    \rowcolor{lightgray}
    \textbf{CNN-LSTM (Focal)} & 0.5581 & 0.2702 & \textbf{0.7553}
      & 0.3980 & 0.7001 \\
    \bottomrule
  \end{tabular}
\end{table}

\noindent The Focal Loss CNN-LSTM caught \textbf{318 out of 421 buggy modules}
in the test set (75.5\%), compared to 190 caught by the best baseline
(Random Forest + SMOTE), representing an improvement of \textbf{128
additional bugs detected} -- a 67.4\% increase in Recall.

% ═══════════════════════════════════════════════════════════
\section{Artefacts Generated}

\begin{table}[H]
  \centering
  \caption{Output Files Generated Today}
  \begin{tabular}{ll}
    \toprule
    \textbf{File} & \textbf{Description} \\
    \midrule
    \texttt{data/raw/jm1.csv}                & Original JM1 dataset \\
    \texttt{data/processed/jm1\_labeled.csv} & Cleaned, binary-labelled dataset \\
    \texttt{data/splits/X\_train.csv}        & Raw training features \\
    \texttt{data/splits/X\_train\_smote.csv} & SMOTE-balanced training features \\
    \texttt{data/splits/X\_test.csv}         & Test features \\
    \texttt{results/01\_class\_distribution.png} & Class imbalance chart \\
    \texttt{results/02\_feature\_distributions.png} & Feature histograms \\
    \texttt{results/03\_feature\_importance\_ratio.png} & Feature ratios \\
    \texttt{results/04\_correlation\_heatmap.png} & Correlation matrix \\
    \texttt{results/05\_smote\_effect.png}   & SMOTE before/after \\
    \texttt{results/11\_baseline\_comparison.png} & Baseline bar chart \\
    \texttt{results/15\_all\_models\_comparison.png} & Full comparison chart \\
    \texttt{models/cnn\_lstm\_bce.pth}       & Saved CNN-LSTM (BCE) weights \\
    \texttt{models/cnn\_lstm\_focal.pth}     & Saved CNN-LSTM (Focal) weights \\
    \texttt{results/all\_results\_so\_far.csv} & All model metrics \\
    \bottomrule
  \end{tabular}
\end{table}

% ═══════════════════════════════════════════════════════════
\section{Next Steps}

\begin{enumerate}[leftmargin=1.5cm]
  \item \textbf{Notebook 05 -- BERT Model:} Implement a BERT-based
    classifier for textual bug report classification using the
    Bugzilla/GitBugs dataset.
  \item \textbf{Notebook 06 -- Enhanced CNN-LSTM:} Fine-tune the
    Focal Loss model with additional epochs and class-weighted
    training to improve F1-Score while maintaining high Recall.
  \item \textbf{10-Fold Cross-Validation:} Apply stratified k-fold
    cross-validation across all models for statistical reliability.
  \item \textbf{Cross-Project Validation:} Train on JM1, evaluate
    on KC1 to assess model generalisability.
  \item \textbf{Report Writing:} Begin drafting Methodology and
    Results chapters using findings from today's implementation.
\end{enumerate}

% ═══════════════════════════════════════════════════════════
\section{Conclusion}

Today's implementation session successfully completed the data
preparation phase and established both baseline and initial deep
learning model results. The Focal Loss variant of the CNN-LSTM model
demonstrated a notably higher bug detection rate (Recall: 0.755)
compared to the best traditional ML baseline (Recall: 0.451),
providing initial empirical evidence supporting the research
contribution of imbalance-aware training for software defect
prediction. The foundation is now in place for the BERT model
implementation and final comparative analysis.

\vspace{1cm}
\noindent\rule{\linewidth}{0.5pt}\\[0.3cm]
\begin{tabular}{ll}
  \textbf{Student:}    & S.Shanthosh (20APSE4882) \\
  \textbf{Supervisor:} & Mrs.\ WMLS Abeythunga \\
  \textbf{Date:}       & 27th May 2026 \\
\end{tabular}

\end{document}
